\begin{pastexam}[label={px:df3-101-ds-05}]{地方三等資訊處理 101 年 第五題}
（申論；20 分）

鏈結串列（linked list）與陣列（array）是常用的資料結構，卻都有其優缺點，使用者往往要依據不同的應用選擇
鏈結串列或陣列來儲存資料以盡全功。請就下列應用說明您會採用的資料結構與原因：（每小題 5 分，共 20 分）
\begin{itemize}[leftmargin=2.2em]
\item[㈠] 二元搜尋（binary search）。
\item[㈡] 要存的資料量很大並懷疑主記憶體沒有一塊連續的空間足以存放這些資料。
\item[㈢] 某應用其資料儲存後經常要增加（insert）或刪除（delete）資料，刪增動作頻繁且資料量不固定。
\item[㈣] 某應用其資料儲存後相對穩定且資料量固定，卻要頻繁尋找或取出資料內容。
\end{itemize}
\end{pastexam}

\begin{pastsoln}
答案：㈠ 陣列；㈡ 鏈結串列；㈢ 鏈結串列；㈣ 陣列。

判斷的依據是兩者的基本差別（資結書第~4~章、第~5.1~節）：
\begin{center}
\small
\begin{tabular}{lll}
\toprule
 & 陣列 & 鏈結串列 \\
\midrule
記憶體 & 一整塊連續空間 & 節點分散，用指標串起來 \\
存取第 $i$ 個 & $O(1)$（位址直接算出來） & $O(i)$（從頭一個一個走） \\
已知位置的插入、刪除 & $O(n)$（要搬移後面的元素） & $O(1)$（改兩個指標） \\
大小 & 事先固定（或整塊重新配置） & 隨時增減 \\
額外空間 & 無 & 每個節點一個（或兩個）指標 \\
\bottomrule
\end{tabular}
\end{center}

\textbf{㈠ 陣列}。 二元搜尋每一步要直接取出區間正中間 $\lfloor (\mathit{low} + \mathit{high})/2 \rfloor$ 那個
元素來比較，需要 $O(1)$ 的隨機存取，總共 $O(\log n)$ 步、$O(\log n)$ 時間（資結書第~17~章）。 鏈結串列要走到
中間就要 $O(n)$，光是第一步的時間就和循序搜尋一樣，失去二元搜尋的意義。 （資料也必須先排好序。）

\textbf{㈡ 鏈結串列}。 陣列要求一整塊連續的記憶體，空間總量夠但被切碎時就配置不到。 鏈結串列每個節點各自
配置，散在記憶體的哪裡都可以，用指標串起來，只要剩餘的空間總和夠（另加每個節點的指標空間）就放得下。

\textbf{㈢ 鏈結串列}。 在已知位置插入、刪除只要改指標，$O(1)$，不必像陣列那樣把後面的元素整批往後或往前搬
（平均 $O(n)$）；資料量不固定時，節點隨用隨配置、刪了就釋放，不必事先估計最大量，也不會因為陣列滿了而要
整塊重新配置和複製。

\textbf{㈣ 陣列}。 資料量固定，事先配置剛好大小的陣列不會浪費，也不會不夠；資料穩定、很少插入刪除，陣列
搬移元素的缺點用不到。 陣列可以 $O(1)$ 依位置取出資料，排好序後可以用二元搜尋 $O(\log n)$ 尋找；資料連續存放，
快取效果好，也不必多存指標。 鏈結串列取第 $i$ 個要 $O(i)$、尋找只能循序 $O(n)$，不適合頻繁查詢。
\end{pastsoln}
